\subsection{Variance - Feature Selection Method}
\subsubsection{Feature Extraction and Selection}

In this experiment, we first transform the raw text data into dense feature vectors using the \textbf{GloVe} embedding method:
\begin{itemize}
    \item Each document is represented as a dense vector capturing semantic information from word embeddings.
\end{itemize}

After feature extraction, we apply Variance Threshold feature selection to reduce dimensionality and retain the most informative features:
\begin{itemize}
    \item The Variance Threshold method removes features with variance lower than a specified threshold (0.01 in this experiment).
    \item Features with sufficient variability across samples are kept.
    \item This step helps eliminate non-informative features and improve both computational efficiency and model performance.
\end{itemize}

\subsubsection{Model Training Procedure}

For each classifier and feature type, the following training procedure is applied:
\begin{itemize}
    \item \textbf{Hyperparameter tuning:} Best hyperparameters are searched using cross-validation techniques (e.g., Grid Search).
    \item \textbf{Cross-validation:} K-fold cross-validation is used to evaluate model stability and avoid overfitting.
    \item \textbf{Loss plotting:} Training curves (e.g., loss vs. epochs) are generated to visualize the learning dynamics.
    \item \textbf{Model saving:} The best performing model is saved for reproducibility and future inference.
\end{itemize}

\subsubsection{Experimental Setup}

The models trained in this study include:
\begin{itemize}
    \item Decision Tree (DT)
    \item Linear Discriminant Analysis (LDA)
    \item Logistic Regression (LogReg)
    \item Multi-Layer Perceptron (MLP)
    \item Perceptron
    \item Random Forest (RF)
    \item Support Vector Machine (SVM)
    \item XGBoost (XGB)
    \item Gaussian Naive Bayes (GNB)
\end{itemize}

All models use features extracted by \textbf{GloVe}. We refer to models as \texttt{ModelName-glove}, e.g., \texttt{DT-glove}, \texttt{MLP-glove}, etc.


\subsubsection{Evaluation Metrics}

For each experiment, we evaluate models based on:
\begin{itemize}
    \item Accuracy
    \item Precision
    \item Recall
    \item F1 Score
    \item ROC AUC
\end{itemize}

\textbf{The results are summarized in the following table:}

\begin{table}[h!]
\centering
\begin{tabular}{lccccc}
\toprule
\textbf{Model} & \textbf{Accuracy} & \textbf{Precision} & \textbf{Recall} & \textbf{F1 Score} & \textbf{ROC AUC} \\
\midrule
DT-glove & 0.6124 & 0.6215 & 0.6286 & 0.6250 & 0.6528 \\
LDA-glove & 0.6881 & 0.6888 & 0.7115 & 0.6999 & 0.7591 \\
LogReg-glove & 0.7147 & 0.7134 & 0.7434 & 0.7281 & 0.\textbf{\underline{7876}} \\
MLP-glove & 0.6975 & 0.7047 & 0.7029 & 0.7038 & 0.7658 \\
Perceptron-glove & 0.6450 & 0.6615 & 0.6259 & 0.6432 & 0.7023 \\
RF-glove & 0.6901 & 0.6876 & 0.7274 & 0.7069 & 0.7649 \\
SVM-glove & 0.6919 & 0.6946 & 0.7090 & 0.7018 & 0.7596 \\
XG-glove & \textbf{\underline{0.7198}} & \textbf{\underline{0.7197}} & \textbf{\underline{0.7449}} & \textbf{\underline{0.7321}} & \textbf{\underline{7939}} \\
GNB-glove & 0.6221 & 0.6403 & 0.6039 & 0.6216 & 0.6707 \\
\bottomrule
\end{tabular}
\caption{Model performances using GloVe feature extraction and feature selection}
\end{table}

\subsubsection{Discussion}
From the results, it can be observed that:
\begin{itemize}
    \item \textbf{XGBoost-glove} outperformed all other models across almost all evaluation metrics.
    \item Logistic Regression and MLP also showed competitive performance, especially in terms of ROC AUC and F1 Score.
    \item Traditional models such as Decision Tree and Gaussian Naive Bayes performed relatively lower, confirming the advantage of more complex classifiers.
\end{itemize}